\begin{table}[htbp]
\centering\small
\caption{Feature-family test AUC, with neural and oracle references.
A dash denotes an unreported cell; N/A indicates that no generating lag exists.
All-lag rows use
stacked gap-specific pair counts. Near-oracle pair-feature performance
establishes feature sufficiency, not a causal explanation of neural predictions.}
\label{tab:signal_decomposition}
\begin{tabular}{lccccc}
\toprule
Method & OC-Det. & OC-Noisy & GPC & KIP-4 & KIP-6 \\
\midrule
Chance & .500 & .500 & .500 & .500 & .500 \\
Key-count baseline & .774 & .574 & .503 & .500 & .500 \\
Full-count XGBoost & .816 & .584 & .503 & .503 & .500 \\
Best contiguous \(k\)-gram & .819 & .586 & .502 & .504 & .504 \\
Known-lag pair LR & .997 & .671 & N/A & N/A & N/A \\
All-lag pair LR ($L_1$) & .997 & .632 & --- & .500 & .502 \\
All-lag pair XGBoost & --- & .667 & --- & .590 & .501 \\
Best neural model & .999 & .671 & .503 & 1.000 & .505 \\
Oracle AUC & 1.000 & .670 & 1.000 & 1.000 & 1.000 \\
\bottomrule
\end{tabular}
\par\smallskip
\begin{minipage}{\linewidth}
\footnotesize
KIP key counts are constant, giving AUC $0.5$ exactly; a known generating
lag is not applicable. KIP $L_1$ results use 100,000 training and 50,000
test sequences. KIP neural references are the highest recorded test AUCs
across completed ordered runs, not seed averages.
\end{minipage}
\end{table}
